\subsection*{S13. Array of objects}

A field whose value is an array of objects.

\begin{lstlisting}[style=json]
{ "id": 1, "items": [ { "sku": "a", "qty": 2 },
                      { "sku": "b", "qty": 1 } ] }
\end{lstlisting}

\options
\begin{enumerate}[label=\textbf{(\alph*)}]
  \item \textbf{Child table with a back-reference and an ordering index},
        and a presence marker in the parent when the path is partial.
        \selected
  \item \textbf{Child table with a back-reference but no index.} The array
        becomes a set rather than a sequence.
  \item \textbf{Parent holds a list of child identifiers}, keeping the
        direction of S7.
  \item \textbf{A junction table} between parent and child, as for a
        many-to-many relationship.
  \item \textbf{Parent references a collection entity}, with the elements
        referencing the same collection:
        \texttt{root(id, items\_id)}, \texttt{item\_list(id)},
        \texttt{items(id, list\_id, idx, ...)}.
\end{enumerate}

\decision{Child table carrying a foreign key back to the parent and an
ordering index. The parent gains no column when the path is total, and a
presence marker when it is partial.}

\begin{center}
\begin{tabular}{@{}l@{}}
\textbf{Table root} \\
\hline
\texttt{id} \\
\hline
1 \\
\hline
\end{tabular}
\qquad
\begin{tabular}{@{}lllll@{}}
\multicolumn{5}{@{}l}{\textbf{Table items}} \\
\hline
\texttt{id} & \texttt{root\_id} & \texttt{idx} & \texttt{sku} & \texttt{qty} \\
\hline
1 & 1 & 0 & a & 2 \\
2 & 1 & 1 & b & 1 \\
\hline
\end{tabular}
\end{center}

\begin{lstlisting}[style=ml]
module Item : sig
  type id
  type t = { id : id; root_id : Root.id; idx : int;
             sku : string; qty : int }
  val get : db -> id -> t
end

module Root : sig
  type id
  type t = { id : id }
  val get : db -> id -> t
  val items : db -> t -> Item.t list
end
\end{lstlisting}

\paragraph{The presence marker.}
An empty array is a \emph{present} value, so it must not be confused with an
absent field. Both produce zero rows in the child table, and the parent is
therefore where the difference has to be recorded. This is not new machinery:
S4 already makes a path that is sometimes absent partial, and producing an
\texttt{option} at all requires the parent to record presence.

\begin{lstlisting}[style=json]
{ "id": 1, "items": [ { "sku": "a" } ] }
{ "id": 2, "items": [] }
{ "id": 3 }
\end{lstlisting}

\begin{lstlisting}[style=ml]
type t = { id : id; items_present : bool }
val items : db -> t -> Item.t list option

  root 1 -> Some [ {sku = "a"} ]
  root 2 -> Some []                (* present, empty *)
  root 3 -> None                   (* absent *)
\end{lstlisting}

Under S4, \texttt{"items": null} collapses with the absent case and yields
\texttt{None}. The marker also appears as a record field, since the record is
the row (S21.3).

\rationale{The ordering index is required rather than optional. JSON array
order is significant, so without it \texttt{["a","b"]} and \texttt{["b","a"]}
shred identically, reconstruction recovers a multiset rather than a list, and
a query asking for the first element has an answer on the JSON side that the
database cannot reproduce. Option (b) is excluded on exactly that ground.

Option (c) is not relational --- a list-valued column is not a column --- and
would still need an index to order it. Option (e) was considered seriously,
since it is more uniform: parent-holds-reference would then be the single rule
for both nested objects and arrays, instead of the direction flipping with
cardinality. It was declined because, absent any sharing of arrays between
documents, \texttt{item\_list} holds exactly one row per parent row. It
carries no information beyond ``this array exists'', which is what the
presence marker records directly, and it costs an extra join on every access.
Option (e) becomes worthwhile only if arrays are given identity so that equal
arrays can be shared, which would make shredding depend on structural
hashing.}

\limitation{A normalization treating \texttt{[]}, \texttt{null} and absence as
equivalent at array-typed paths was considered. It would remove the presence
marker entirely and make array accessors total, at the cost of stating the
equivalence in the specification and accepting that the resulting theorem is
relative to that semantics rather than to the ones implemented by jq, JSONPath
or SQL/JSON, all of which distinguish \texttt{null} from \texttt{[]}. It is
recorded here as an available simplification, not adopted.

As with nullability generally, the parent table's shape depends on what was
observed: a path that is total in the corpus has no presence marker, and one
that is partial has one.}
